\begin{proof}
Procederemos en varias etapas.

I) \emph{Reducción al caso en que el anillo $B$ es reducido.} ---
Pongamos $B_0=B_{\mathrm{red}}=B/\mathfrak N$, donde $\mathfrak N$
es el nilradical de $B$, y sea $C_0$ la clausura integral de $A$ en
$B_0$. Se tiene el siguiente lema:

\begin{lemma}[6.14.4.1]
\label{IV.2.s6.14.4.1-fr}
Sean $A$ un anillo, $B$ una $A$-álgebra y $B_0=B/\mathfrak n$ el
cociente de $B$ por un ideal cuyos elementos son nilpotentes. Si
$C_0$ es la clausura integral de $A$ en $B_0$, la imagen inversa
$C$ de $C_0$ por el homomorfismo canónico $\varphi:B\to B_0$ es la
clausura integral de $A$ en $B$.
\end{lemma}

